\documentclass[11pt]{article}
\usepackage[a4paper,margin=0.9in]{geometry}
\usepackage{booktabs}
\usepackage{amsmath}
\usepackage{amssymb}
\usepackage{graphicx}
\usepackage{graphics}
\usepackage{hyperref}
\usepackage{microtype}
\hypersetup{colorlinks=true,linkcolor=black,urlcolor=blue}
\emergencystretch=1.2em

\title{Layer Reduction with a Fitted Ternary TerGRASP Bridge}
\author{Sumit Jha}
\date{October 2026}

\begin{document}
\maketitle
\sloppy

\begin{abstract}
We test whether the most redundant decoder layers of a small LLM can be
replaced by a $1.58$-bit ternary low-rank bridge whose parameters are
\emph{fitted} to the layer's residual. We rank layers with Block Influence
(BI), remove/replace the $k$ lowest-BI blocks, and measure validation
perplexity (PPL) on wikitext subsets with a random-removal control. The
fitted ternary bridge matches plain hard removal at $k{=}1$ while keeping
the graph depth and shrinking the replaced blocks to $\sim$1.58 bit/param.
Naive averaging, weight tying, and a single affine bridge across the gap all
fail by orders of magnitude, which is what motivated fitting the bridge per
layer. Supporting runs (SVD low-rank, quantized no-op bypass) confirm the
same BI ranking drives both pruning and quantization choices.
\end{abstract}

\section{Report}
\subsection{Mentor feedback since Presentation 1}
FILL IN. In short: an earlier plan relied on random TerGRASP bridges, which
the mentor flagged because the bridge is then $\approx$ identity. We changed
the method to \emph{fit} the bridge on the residual Delta $= h_{out}-h_{in}$
and added the random-removal control and the per-approach reports listed in
Section~\ref{sec:approaches}.

\subsection{Baselines and related work}
We compared against the state of the art that is runnable on our budget:
\begin{itemize}
\item \textbf{ShortGPT (Men et al.)} - the BI ranking we reproduce as our
control. Targeted removal of low-BI blocks degrades PPL far more slowly
than random removal.
\item \textbf{LayerDrop (Fan et al.)} - fine-tune with random layer drops.
Archived; the 60-step CPU/MPS budget was too small to separate depths.
\item \textbf{CALM / early exits}, \textbf{KD to shallow students}, and
\textbf{L0 gate pruning} - archived for the same reason.
\end{itemize}
No closed-source large-LLM baselines were reproduced, per the instructions.

\subsection{Datasets, bias, bottlenecks}
Main corpus: Salesforce/wikitext, configuration wikitext-2-raw-v1;
comparison configurations wikitext-103-raw-v1 and wikitext-2-v1.
\paragraph{Bias handling.} We reused the Assignment 1 preprocessing
(\texttt{join\_nonempty}, non-overlapping 256-token windows) and capped
\texttt{train} to 300k chars, validation/test to 100k chars, EVAL\_TOKENS
to 3000, seed fixed to 42. This keeps the PPL windows identical across
approaches, so no method gets an easier slice.
\paragraph{Bottleneck.} We initially wanted PTB and tiny\_shakespeare, but
datasets 5.x removed the script loaders and their parquet mirrors were not
resolvable for those identifiers, so we rotated the four wikitext
configurations instead. The principal bottleneck otherwise was the single
MPS laptop: training-heavy approaches (KD, LayerDrop, gates) never got a
large enough budget, which is why their archived rows are non-separating.
\paragraph{Dataset extension.} None was introduced; the differences from
previous work are in the preprocessing and the 3k-token PPL window, which
we declare explicitly.

\subsection{Experimental plan}
\paragraph{Models.} GPT-2 small (124,439,808 params, 12 blocks) and
SmolLM-135M (134,515,008 params, 30 blocks). Architectural accessors
differ: GPT-2 stores blocks in \texttt{model.transformer.h}, SmolLM in
\texttt{model.model.layers}; a thin helper (\texttt{get\_layers},
\texttt{set\_layers}) normalizes both.
\paragraph{How approaches differ.} (i) \emph{target selection} is always
BI-ranked vs random; (ii) the \emph{replacement operator} differs: none
(hard removal), fitted ternary bridge (TerGRASP), rank-$r$ SVD truncation,
2/4/8-bit quantization; (iii) \emph{scope} differs: one layer ($k{=}1$) up
to four ($k{=}4$).
\paragraph{Hyperparameters and selection.} Bridge rank $r{=}32$; bridge
optimizer Adam, lr $10^{-2}$, 80 steps; AbsMean ternary quantization;
ridge $10^{-4}$. Ablations: we swept tokenization budgets, EVAL\_TOKENS,
random seeds for the control, bridge rank $\in\{64,128,256\}$ via the SVD
variant, and bitwidth $\in\{8,4,2\}$. Selection of layers is always the BI
ranking, never a tuned heuristic.
\paragraph{Metric.} Primary: validation perplexity,
$\mathrm{PPL}=\exp(\frac{1}{N}\sum_i -\log p_\theta(x_i|x_{<i}))$ over
256-token non-overlapping windows, identical for all approaches.
Secondary: latency in ms/forward pass and parameter count. We also
recorded the BI value of each removed layer to check that low-BI removal
is indeed the cheapest.
\paragraph{System effort.} One MacBook, MPS backend for GPT-2 and SmolLM,
CPU for the 02 pass. A full 15-approach run is ~1 hour; the final
\texttt{assignment2.ipynb} re-runs in a few minutes and embeds the final
tables and the two figures.

\subsection{Project management}
\paragraph{Novelty proposal.} Replace the random TerGRASP bridge with a
fitted low-rank bridge and ternarize its factors; the BI ranking from
Assignment 1 then chooses \emph{which} block to replace.
\paragraph{Resources.} 1 machine, 16 GB RAM, Apple MPS, no paid GPU.
\paragraph{Split.} FILL IN per-member: dataset/eval code, TerGRASP
implementation, SVD+quantization supporting runs, figures/report.

\section{Experiment}
\subsection{Implementation}
\texttt{assignment2.ipynb} runs end-to-end on MPS: dataset prep, BI
ranking, residual capture via forward hooks, bridge fitting, ternary
quantization, block replacement, PPL/latency/parameter eval, two figures.
\subsection{Approaches tried (all 15)} \label{sec:approaches}
\small
\small
\resizebox{\textwidth}{!}{\begin{tabular}{llp{6.5cm}}
\toprule
\# & Approach & Outcome \\
\midrule
01 & Fitted-ternary TerGRASP & headline result, works \\
02 & Affine bridge across gap & fails: PPL $10^3$-$10^8\times$ \\
03 & BI + KD repair & scaffold, KD loss $\approx 0$ at 60-step budget \\
04 & Permutation-aligned averaging & fell back to naive averaging, fails \\
05 & LayerDrop->BI->SVD->KD pipeline & scaffold, identical rows \\
06 & ShortGPT BI baseline & works, the control \\
07 & SVD low-rank & works on SmolLM; GPT-2 row needs clean rerun \\
08 & Weight-space merging & always worse PPL \\
09 & KD shallow student & partial recovery at tiny budget \\
10 & L0 learned gates & gates never close, scaffold \\
11 & LayerDrop & flat PPL, scaffold \\
12 & Attn/FFN proxy pruning & stub \\
13 & Early exit & scaffold, no exit-head training \\
14 & Cross-layer weight tying & PPL $10^3$-$10^6\times$, fails \\
15 & Quantized no-op bypass & works, supporting rule \\
\bottomrule
\end{tabular}}
\normalsize

\subsection{Results}
\begin{table}[h]
\centering
\small
\small
\resizebox{\textwidth}{!}{\begin{tabular}{lrrrrr}
\toprule
Model & k=1 & k=2 & k=3 & k=4 & Note \\
\midrule
GPT-2 fitted (this work) & 45.11 & 50.09 & 140.42 & 879.52 & keeps graph depth, $\sim$1.58-bit blocks \\
GPT-2 hard removal & 45.76 & 47.95 & 121.81 & 943.24 & control \\
SmolLM fitted (this work) & 21.71 & 40.49 & 33.05 & 38.49 & \\
SmolLM hard removal & 19.93 & 26.13 & 29.91 & 34.20 & \\
\bottomrule
\end{tabular}}
\normalsize
\caption{Validation PPL by number of replaced layers on the wikitext-2-raw-v1 3k-token window.}
\end{table}
\medskip

For Approach 06, the targeted/random control at the same k (PPL):
\small
\small
\resizebox{\textwidth}{!}{\begin{tabular}{lrrrrrr}
\toprule
model & \multicolumn{4}{c}{BI targeted (k=1..4)} & random k=1 & random k=3 \\
\midrule
GPT-2   & 45.76 & 47.95 & 121.81 & 943.24 & 44.42 & 4181.21 \\
SmolLM  & 19.93 & 26.13 & 29.91  & 34.20  & 20.83 & 34.23   \\
\bottomrule
\end{tabular}}
\normalsize

\subsection{Findings}
At $k{=}1$ the fitted ternary bridge matches or slightly beats hard removal
while keeping the graph depth and cutting that block's storage to $\sim$1.58
bit/param. For $k\ge 3$ it degrades in line with hard removal; it does not
blow up like averaging/tying, and it does not pretend to add a real layer
like the random bridge. The negatives (02, 08, 14) show that the layer maps
are neither affine nor weight-interchangeable, which is exactly why a
\emph{fitted} low-rank bridge is the minimal fix.

\subsection{What ran successfully and why}
Only four notebooks produced a usable signal under the free MPS budget:
01 (fitted ternary, headline), 06 (BI control), 07 (SmolLM SVD curve), 15
(BI-guided 2-bit quantization). The rest were archived: their frameworks
ran but the signal was flat, or they failed by orders of magnitude.

\section{Conclusion}
The same BI ranking that makes layer removal cheap also predicts that the
kept matrix of a low-BI block can be ternarized or SVD-truncated with no
PPL loss at $k{=}1$. A fitted ternary low-rank bridge inserted at that
position preserves validation PPL on both GPT-2 and SmolLM-135M, keeps the
graph depth, and shrinks the replaced blocks to $\sim$1.58 bit/param. This
is the approach we recommend for the assignment.

\bibliographystyle{plain}
\nocite{*}
\begin{thebibliography}{9}
\bibitem{shortgpt}
Men et al., \emph{ShortGPT: Layers in Large Language Models are More
Redundant Than You Expect}, 2024.
\bibitem{layerdrop}
Fan et al., \emph{Reducing Transformer Depth on Demand with Structured
Dropout} (LayerDrop), ICLR 2020.
\bibitem{calm}
Schwartz et al., \emph{The Right Tool for the Job: Matching Model and
Instance Complexities} (CALM), 2020.
\bibitem{bitnet}
Ma et al., \emph{The Era of 1-bit LLMs: All Large Language Models are in
1.58 Bits}, 2024.
\bibitem{kd}
Hinton et al., \emph{Distilling the Knowledge in a Neural Network}, 2015.
\end{thebibliography}

\section{Appendix: Why the Fitted Bridge Works (maths)}
Let $x$ be the hidden state entering layer $L$, and suppose the block is
close to identity: $F_L(x) = x + \Delta(x)$ with
$\|\Delta(x)\| \ll \|x\|$. Hard removal replaces $F_L$ by the identity
$g(x)=x$, incurring error $\Delta(x)$. We propose $g(x) = x + xUV$ and
choose $U,V$ to solve
\[
\min_{U,V} \; \bigl\| XUV - \Delta \bigr\|_F^2,
\quad \Delta_i = F_L(x_i) - x_i .
\]
For a near-identity layer $\Delta$ is small and lives in a low-dimensional
cone (the directions in which the residual actually moves), so a rank-$r$
factor model recovers most of its energy: by the Eckart-Young theorem the
best rank-$r$ approximation to the centered residual has error equal to the
tail $\sum_{i>r}\sigma_i^2$, which is small precisely when $\Delta$ is
low-rank. Ternary quantization then perturbs the factors by
$\|W - W_q\|_\infty \le \gamma$, so the bridge output moves by at most
$\mathcal{O}(\gamma^2 \|x\|)$. Putting the two together,
\[
\| g(x) - F_L(x) \| \;\le\; \underbrace{\sum_{i>r}\sigma_i}_{\text{SVD tail}}
\;+\; \underbrace{\mathcal{O}(\gamma^2 \|x\|)}_{\text{ternary noise}},
\]
which is small when the layer is low-BI. This is exactly the regime the
BI ranking selects, so the choice of *which* layer to replace and the form
of the replacement are the same decision.

\section{Appendix: Per-step result tables}
\subsection{Step tables}
\small
\small
\resizebox{\textwidth}{!}{\begin{tabular}{llll}
\toprule
Step & GPT-2 value & SmolLM value & Meaning \\
\midrule
Model params & 124,439,808 & 134,515,008 & nearly equal \\
\#blocks & 12 & 30 & \\
k=1 fitted PPL & 45.11 & 21.71 & one block replaced \\
k=1 removal PPL & 45.76 & 19.93 & \\
rank-256 SVD PPL & 36.78 (guard) & 19.93 & full-rank \\
2-bit on L5 PPL & 36.78* & 17.76 & L5 tolerated \\
\bottomrule
\end{tabular}}
\smallskip
\normalsize
*GPT-2 row is the dtype-guard artifact discussed in Section~\ref{sec:approaches}.

\section{Conclusion (comparison study)}
\small
\small
\resizebox{\textwidth}{!}{\begin{tabular}{lllll}
\toprule
Approach & Works? & PPL effect & Compression? & Role \\
\midrule
01 Fitted-ternary TerGRASP & yes & $\approx$ baseline at k=1 & $\sim$20$\times$ on replaced block & headline \\
06 ShortGPT BI & yes & targeted $\gg$ random at k$\ge$2 & none & control \\
07 SVD low-rank & yes & smooth rank/PPL trade & 2-4$\times$ on shrink blocks & support \\
15 2-bit no-op bypass & yes & $\approx$ no change on low-BI & up to $16\times$ on block & rule of thumb \\
02/08/14 merge/tying/bridge & no & PPL 10$^3$-10$^6\times$ & irrelevant & negative results \\
03/09/10/11/12/13 & scaffold & no separation & - & future work \\
\bottomrule
\end{tabular}}
\normalsize

Among all 15, \textbf{01 is the best} because it is the only one that (a)
keeps the graph depth, (b) shrinks the replaced block to $\sim$1.58
bit/param, and (c) preserves PPL at the operating point $k{=}1$ on both
models. 06 remains as the control, 07 and 15 are supporting evidence that BI
predicts compressibility, and 02/08/14 are the negative results that
motivate fitting rather than averaging or tying.

\end{document}
